\section*{Capítulo \ref{ch:b2:retrosynthesis} --- Retrossíntese e estratégia em várias etapas}

\begin{solution}{exo:b2:retrosynthesis:1}
\begin{itemize}
  \item 2-fenilpropan-2-ol: propanona + \ce{PhMgBr}, ou acetofenona + \ce{CH3MgBr}.
  \item Pentan-3-ol: propanal + \ce{CH3CH2MgBr}.
  \item 1-ciclo-hexiletanol: etanal + brometo de ciclo-hexilmagnésio, ou ciclo-hexanocarbaldeído +
        \ce{CH3MgBr}.
  \item 2-metilbutan-2-ol: propanona + \ce{CH3CH2MgBr}, ou butanona + \ce{CH3MgBr}.
  \item 1-feniletanol: benzaldeído + \ce{CH3MgBr}, ou etanal + \ce{PhMgBr}.
\end{itemize}
\end{solution}

\begin{solution}{exo:b2:retrosynthesis:2}
\omterm{def:b2:retrosynthesis:disconnection}{Sínton} aceptor $\mathrm{Ph\overset{+}{C}H{-}OH}$, equivalente o benzaldeído; \omterm{def:b2:retrosynthesis:disconnection}{sínton} doador \ce{CH3-},
equivalente o brometo de metilmagnésio (ou o metil-lítio). (O outro corte, com \ce{Ph-} como doador,
usa \ce{PhMgBr} e o etanal.)
\end{solution}

\begin{solution}{exo:b2:retrosynthesis:3}
$0.90 \times 0.85 \times 0.80 \times 0.95 \times 0.70 = 0.407$: cerca de 41~\%.
\end{solution}

\begin{solution}{exo:b2:retrosynthesis:4}
O \ce{LiAlH4}, necessário para reduzir o éster ao álcool, reduziria também a cetona. Proteja a
cetona na forma de seu acetal cíclico (etano-1,2-diol, catalisador ácido, água removida); reduza o
éster com \ce{LiAlH4}; remova o acetal com ácido aquoso diluído: 5-hidroxipentan-2-ona.
\end{solution}

\begin{solution}{exo:b2:retrosynthesis:5}
IGF: butano-1,3-diol $\Rightarrow$ 3-hidroxibutanal (redução do aldeído, \ce{NaBH4}). \omterm{def:b2:retrosynthesis:disconnection}{Desconexão} a
dois grupos desse $\beta$-hidroxialdeído na ligação $\alpha$--$\beta$: duas moléculas de etanal
(\omterm{def:b2:enolates-aldol:aldol}{aldólica}). No sentido direto: etanal, base diluída, a frio; depois \ce{NaBH4}.
\end{solution}

\begin{solution}{exo:b2:retrosynthesis:6}
O anel de ciclo-hexeno é cortado nas duas ligações \ce{C-C} criadas em uma \omterm{def:b2:frontier-orbitals:diels-alder}{reação de Diels--Alder},
as alílicas em relação à ligação dupla do anel, do lado do substituinte: buta-1,3-dieno e
but-3-en-2-ona, com o grupo acetila no \omterm{def:b2:frontier-orbitals:diels-alder}{dienófilo}.
\end{solution}

\begin{solution}{exo:b2:retrosynthesis:7}
Numerando C1 (\ce{C=O}), C2=C3, depois C4 com as duas metilas
(\cref{met:b2:conjugate-additions:robinson-disconnection}): o corte aldólico C2--C3 deixa um aldeído
em C3 e uma metilcetona; o corte de Michael C4--C5 deixa o 2-metilpropanal (C3 sua carbonila, C4 seu
carbono $\alpha$) e a but-3-en-2-ona.
\end{solution}

\begin{solution}{exo:b2:retrosynthesis:8}
Uma etapa: benzeno com cloreto de butanoíla e \ce{AlCl3} (\omterm{def:b2:aromatic-substitution:friedel-crafts}{acilação de Friedel--Crafts}, sem rearranjo,
acilação única). Duas etapas: benzaldeído com brometo de propilmagnésio, depois oxidação do álcool
secundário. A via de Friedel--Crafts é mais curta e seus reagentes mais baratos, mas usa um equivalente
inteiro de cloreto de alumínio, que termina nos efluentes aquosos.
\end{solution}

\begin{solution}{exo:b2:retrosynthesis:9}
O clorocromato de piridínio no diclorometano, a oxidação de Swern ou o periodinano de Dess--Martin
dão o hexanal; o dicromato aquoso acidificado ou o permanganato dão o ácido hexanoico, porque em água
o aldeído forma seu hidrato, que é oxidado de novo (\cref{prop:b2:retrosynthesis:mild-oxidation}).
\end{solution}

\begin{solution}{exo:b2:retrosynthesis:10}
A amina reage primeiro, sem proteção, se o álcool estiver protegido: proteja o álcool como éter
silílico (ele sobrevive à acilação); acile a amina; remova o grupo silila com fluoreto; acile o
álcool. Se a amina tivesse de ser mantida livre enquanto o álcool reage, ela seria protegida como
carbamato de \textit{terc}-butoxicarbonila (removido por ácido), ortogonal ao éter silílico (removido
pelo fluoreto): qualquer um dos dois pode ser liberado sem tocar no outro.
\end{solution}

\begin{solution}{exo:b2:retrosynthesis:11}
\ce{CH3COCH2CH2COCH3} é uma 1,4-dicarbonila: corte a ligação \ce{C-C} central. O doador é o \omterm{def:b2:enolates-aldol:enolate}{enolato}
do 3-oxobutanoato de etila (etóxido de sódio); o aceptor, um carbono $\alpha$ de cetona, é fornecido
pela 1-bromopropan-2-ona, \ce{BrCH2COCH3} (substituição bimolecular do brometo). Depois a hidrólise
aquosa do éster e o aquecimento removem o \ce{CO2}: hexano-2,5-diona.
\end{solution}

\begin{solution}{exo:b2:retrosynthesis:12}
Retrossíntese: a metilcetona com um \ce{CH2} vizinho é o produto da síntese acetoacética; corte a
ligação entre C3 e C4: o \omterm{def:b2:enolates-aldol:enolate}{enolato} do 3-oxobutanoato de etila e o haleto alílico
1-bromo-3-metilbut-2-eno, \ce{(CH3)2C=CHCH2Br}. No sentido direto: etóxido de sódio em etanol, depois o
haleto; depois base aquosa, acidificação e aquecimento, que hidrolisam o éster e removem o \ce{CO2}
(\cref{prop:b2:enolates-aldol:decarboxylation}): 6-metil-hept-5-en-2-ona.
\end{solution}

\begin{solution}{pb:b2:retrosynthesis:1}
\textbf{1.} \ce{HO-C6H4-CH2CH2COCH3}, o \ce{OH} em para da cadeia: um fenol e uma cetona.
\textbf{2.} Hidrogenação de uma \ce{C=C}: \ce{ArCH2CH2CO} $\Rightarrow$ \ce{ArCH=CHCO}, uma
cetona $\alpha,\beta$-insaturada.
\textbf{3.} (\textit{E})-4-(4-hidroxifenil)but-3-en-2-ona, \ce{HOC6H4CH=CHCOCH3}.
\textbf{4.} Uma carbonila $\alpha,\beta$-insaturada: uma \omterm{def:b2:enolates-aldol:aldol}{condensação aldólica}, cortada na \ce{C=C}.
\textbf{5.} O 4-hidroxibenzaldeído e a propanona.
\textbf{6.} O aldeído não tem hidrogênio $\alpha$ e só pode ser o eletrófilo; a propanona, em excesso,
é a única fonte de \omterm{def:b2:enolates-aldol:enolate}{enolato} e reage com o aldeído, mais reativo, em vez de consigo mesma.
\textbf{7.} (1) 4-hidroxibenzaldeído, propanona em excesso, hidróxido de sódio aquoso; acidificação. (2)
\ce{H2} sobre paládio em carvão, à temperatura ambiente, até a absorção de um equivalente.
\textbf{8.} O grupo \ce{CHO} estabiliza o fenóxido, de modo que esse fenol é pelo menos tão ácido quanto
o fenol ($\mathrm pK_a$ 9.99). Em $\mathrm{pH} > 13$, $[\ce{ArO-}]/[\ce{ArOH}] > 10^{13-10} = 10^3$: o
fenóxido.
\textbf{9.} O \ce{O-} empurra densidade eletrônica através do anel até o carbono carbonílico (um doador
por ressonância em para dele): o aldeído é menos eletrofílico, e a condensação mais lenta.
\textbf{10.} Um éter benzílico: a hidrogenação sobre paládio que reduz a \ce{C=C} também cliva o éter
benzílico (hidrogenólise), liberando o fenol na mesma etapa.
\textbf{11.} Três: benzilação, \omterm{def:b2:enolates-aldol:aldol}{condensação aldólica}, hidrogenação com desproteção.
\textbf{12.} Não: em condições brandas (temperatura ambiente, paládio), o anel \omterm{def:b2:huckel:aromatic}{aromático} não é reduzido.
\textbf{13.} Não proteger: o fenóxido apenas retarda a condensação, o que um tempo maior ou um leve
aquecimento compensam; uma etapa economizada vale mais que o ganho de velocidade.
\textbf{14.} O 4-iodofenol e a but-3-en-2-ona.
\textbf{15.} \ce{HOC6H4I + CH2=CHCOCH3 + Et3N -> HOC6H4CH=CHCOCH3 + Et3NH+ + I-}, catalisador de paládio.
\textbf{16.} No \ce{CH2} terminal, o carbono $\beta$: o grupo arila se insere na extremidade menos
impedida, e a eliminação de hidreto $\beta$ dá então a enona (\textit{E}) conjugada.
\textbf{17.} \omterm{def:b2:catalytic-cycles:oxidative-addition}{Adição oxidativa} do iodeto de arila ao \ce{Pd(0)}, coordenação e inserção do alceno,
eliminação de hidreto $\beta$, e regeneração do \ce{Pd(0)} pela base, que capta o \ce{HI}
(\cref{ch:b2:catalytic-cycles}).
\textbf{18.} $162.188/(220.005 + 70.091 + 101.193) = 162.188/391.289 \approx 41.4~\%$.
\textbf{19.} \ce{C7H6O2 + C3H6O -> C10H10O2 + H2O}: $162.188/(122.123 + 58.080) = 162.188/180.203
\approx 90.0~\%$.
\textbf{20.} 100~\%: todos os átomos do \ce{H2} e da enona terminam no produto.
\textbf{21.} $164.204/(122.123 + 58.080 + 2.016) = 164.204/182.219 \approx 90.1~\%$.
\textbf{22.} A \ce{C=C}: conjugada e desimpedida, ela se liga ao paládio e é hidrogenada muito mais
depressa que a \ce{C=O} da cetona; parar após um equivalente de hidrogênio conserva a cetona.
\textbf{23.} Via \omterm{def:b2:enolates-aldol:aldol}{aldólica}: aldeído, propanona e hidróxido de sódio baratos, água como subproduto. Via
de Heck: um iodeto de arila, um catalisador de paládio, e a but-3-en-2-ona, muito tóxica.
\textbf{24.} A via \omterm{def:b2:enolates-aldol:aldol}{aldólica}: duas etapas, reagentes baratos e mais seguros, economia atômica acima de 90~\%.
\textbf{25.} Ela o multiplica por 0.90: $0.782 \times 0.90 \approx 70~\%$.
\textbf{26.} $0.85 \times 0.92 = \boldsymbol{\approx 78~\%}$ no total.
\end{solution}
